% =============================================================================
%  SLIDE 6 – Connecting to the New Lesson: angles adjacent to a side
%    6a (2 frames): 1 title only   2 the big question
%    6b (6 frames): 1 title only   2 question (left)   3 triangle ABC (right)
%                   4 side AB + brace   5 ∠A (adjacent)   6 ∠B (adjacent)
%  Angle notation: prefix symbol \angle A, \angle B, \angle C (no hats).
% =============================================================================

% ---- 6a : big question only (no diagram) ------------------------------------
\begin{frame}
\begin{tikzpicture}[remember picture,overlay,shift={(current page.south west)},x=1cm,y=1cm]
  \fill[white] (0,0) rectangle (16,9);
  \node[anchor=west,text=NavyDeep] at (0.8,8.35)
       {{\fontsize{12}{17}\selectfont\muol{ទំនាក់ទំនងមេរៀនថ្មី}}};
  \draw[NavyDeep!18,line width=0.6pt] (0.8,7.98)--(15.2,7.98);

  \node<2->[inner sep=0pt,text width=13.6cm,align=center,text=NavyDeep] at (8,4.4)
       {\fontsize{22}{36}\selectfont\bfseries
        តើយើងត្រូវស្គាល់ធាតុប៉ុន្មាននៃត្រីកោណពីរ\\[0.3cm]
        ទើបអាចសន្និដ្ឋានបានថា\\[0.3cm]
        ត្រីកោណទាំងពីរនោះប៉ុនគ្នា?};

  \fill[NavyDeep] (0,0) rectangle (16,0.65);
  \shade[shading=axis,left color=GoldAccent,right color=GoldAccent!0] (0,0.65) rectangle (16,0.71);
  \node[anchor=west,inner sep=0,text=white] at (0.8,0.28)
       {\fontsize{8.5}{11}\selectfont \TeacherName\ \ |\ \ \color{SkyLight}គណិតវិទ្យា ថ្នាក់ទី៨};
  \node[anchor=east,inner sep=0,text=GoldAccent] at (15.2,0.28)
       {\fontsize{8.5}{11}\selectfont\muol{មេរៀនទី១២}\ \ \color{white}ការប្រៀបធៀបត្រីកោណ};
\end{tikzpicture}
\end{frame}
